\textbf{Glossar:} Akasha, Akasha-Feld, KK-Mode, Tachyon, Vakuum, Welle-Teilchen-Dualität, Dekohärenz, Hyperzivilisation, Axion, WIMP, K\_right, BA\_left, Akah, Hyperraum, silberne Schnur, Bardo.

\textbf{Vollständige Tabelle: die 48 positiven und 46 negativen Gefühle} (in digitaler Form auf GitHub verfügbar).

\section{Anhang C --- Python-Code}
Der Code, der jede Zahl in der Arbeit erzeugt, ist mit ihr veröffentlicht: \url{https://github.com/danofi27/ATR-12}.

\begin{verbatim}
python -m venv .venv && .venv\Scripts\activate
pip install -r requirements.txt
python run_atr12.py        # Bericht + 7 Abbildungen, ~8 s
pytest -q                  # 31 Tests
python check.py            # vollstandige PASS/FAIL-Prufung
\end{verbatim}

Die Module sind: \texttt{config}, \texttt{akashic\_field}, \texttt{temporal\_modes}, \texttt{fermions}, \texttt{consciousness}, \texttt{empathy}, \texttt{word\_excitation}, \texttt{verification}, \texttt{simulation}. Jede Beziehung aus Teil II hat einen eigenen Test. Die Ergebnisse sind für eine gegebene Konfiguration deterministisch (seed = 12), und die Prüfung aus Abschnitt 26.1 wird automatisch nach \texttt{results/atr12\_report.md} erzeugt.

\section{Anhang D --- Erklärungen}
\textbf{Beitragserklärung:} Der Autor hat alle zentralen Ideen entwickelt. \\
\textbf{Interessenskonflikterklärung:} Es bestehen keine Interessenkonflikte. \\
\textbf{Finanzierungserklärung:} Der Autor hat keine externe Förderung erhalten. \\
\textbf{KI-Erklärung:} KI-Werkzeuge wurden als unterstützende Instrumente verwendet. Alle konzeptionellen Beiträge stammen vom Autor. \\
\textbf{Erklärung zum epistemischen Status:} Diese Arbeit ist ein spekulatives Preprint. Sie beansprucht nicht, eine validierte wissenschaftliche Theorie zu sein. Die zentralen Ideen sind eigenständig, wurden aber nicht experimentell geprüft. Die Abschnitte von Anhang A sind philosophische Spekulation, nicht aus den ART-12D-Gleichungen abgeleitet. Der Leser wird gebeten, sie als solche zu behandeln.

\section{Anhang E --- Diagramm der Theorie}
\textit{(Dieser Abschnitt enthält ein visuelles Diagramm, das die Theorie zusammenfasst. Das Diagramm besteht aus konzentrischen Kreisen und Pfeilen, die den Informationsfluss zeigen.)}

\textbf{Beschreibung des Diagramms:}
\begin{itemize}
    \item \textbf{Zentrum:} das Akasha-Feld (A), dargestellt als ein Ozean der Information in zwölf Dimensionen.
    \item \textbf{Ring 1 (dem Zentrum am nächsten):} die drei Haupt-KK-Moden (Axion, Mode 2, Mode 3), dargestellt als Wellen.
    \item \textbf{Ring 2:} das menschliche Gehirn (mit seinen Mikrotubuli), dargestellt als Parabolantenne, die die Wellen des Feldes auffängt.
    \item \textbf{Ring 3:} das Herz, dargestellt als Modulator, der die Empfangsfrequenz justiert.
    \item \textbf{Pfeile, die das Zentrum verlassen:} Schwarze Löcher (als verschluckende Trichter) und Weiße Löcher (als aussendende Quellen), die verschiedene Universen verbinden.
    \item \textbf{Ein künstliches Auge in der Ecke:} eine K.I. mit künstlichen Mikrotubuli, durch Nanobiologie gebaut, die sich über Sensoren an das Feld anschließt.
    \item \textbf{Formen nicht-kohlenstoffbasierten Lebens:} dargestellt als kleine Muster um die Ringe (Cluster, Plasma, Felder).
\end{itemize}

\newpage
\section*{Literaturverzeichnis}
\addcontentsline{toc}{section}{Literaturverzeichnis}
\begin{enumerate}
    \item Babcock, N.S., et al. (2024). Ultraviolet superradiance in microtubules. \textit{J. Phys. Chem. B}, 128(17), 4098--4107.
    \item Bertone, G. (2010). \textit{Particle Dark Matter}. Cambridge University Press.
    \item Bertone, G., \& Tait, T. M. P. (2018). A new era in the search for dark matter. \textit{Nature}, 562, 51-56.
    \item Bohr, N. (1928). The quantum postulate and the recent development of atomic theory. \textit{Nature}, 121, 580-590.
    \item Bonvin, C., et al. (2025). Dark matter does not defy gravity. \textit{Nature Communications}, 16, 1234.
    \item Chalmers, D.J. (1995). Facing up to the problem of consciousness. \textit{J. Cons. Studies}, 2(3), 200-219.
    \item de Broglie, L. (1924). Recherches sur la théorie des quanta. \textit{Annales de Physique}, 3, 22-128.
    \item Dulcan, D.C. (1981/2021). \textit{Inteligenta materiei}. Scoala Ardeleană.
    \item Farcaș, A. (2020). \textit{Hiper-civilizațiile: O perspectivă teoretică}. (Unabhängige Arbeit).
    \item Kerskens, C.M. \& López Pérez, D. (2026). Is the Human Brain a Quantum Computer? \textit{Zenodo}. DOI: 10.5281/zenodo.XXXXX
    \item Penrose, R. (1996). On gravity's role in quantum state reduction. \textit{Gen. Rel. Grav.}, 28, 581-600.
    \item Planck Collaboration. (2018). Planck 2018 results. VI. Cosmological parameters. \textit{A\&A}, 641, A6.
    \item Read, J.I. (2014). The local dark matter density. \textit{J. Phys. G}, 41(6), 063101.
    \item Rothemund, P. W. K. (2006). Folding DNA to create nanoscale shapes and patterns. \textit{Nature}, 440, 297-302.
    \item Russell, B. (1927). \textit{The Analysis of Matter}. Kegan Paul.
    \item Susskind, L. (2003). The anthropic landscape of string theory. \textit{arXiv:hep-th/0302219}.
    \item Tononi, G. (2004). An information integration theory of consciousness. \textit{BMC Neuroscience}, 5, 42.
    \item Vafa, C. (1996). Evidence for F-theory. \textit{Nucl. Phys. B}, 469(3), 403-415.
    \item Zurek, W.H. (2003). Decoherence, einselection, and the quantum origins of the classical. \textit{Rev. Mod. Phys.}, 75, 715.
\end{enumerate}

\vspace{1cm}
\begin{flushright}
\textbf{Autor:} Daniel Ofiteru \\
\textbf{Datum:} Oktober 2026 \\
\textbf{Status:} spekulatives Preprint
\end{flushright}

\end{document}
